\documentclass[11pt]{article}
\usepackage[T1]{fontenc}
\usepackage{lmodern,amsmath,amssymb,amsthm,mathtools}
\usepackage[margin=1.05in]{geometry}
\usepackage{longtable,array}
\usepackage[hidelinks]{hyperref}
\newtheorem{theorem}{Theorem}[section]
\newtheorem{lemma}[theorem]{Lemma}
\newtheorem{proposition}[theorem]{Proposition}
\newtheorem{corollary}[theorem]{Corollary}
\theoremstyle{remark}\newtheorem{remark}[theorem]{Remark}
\newcommand{\R}{\mathbb R}
\newcommand{\N}{\mathbb N}
\newcommand{\norm}[1]{\left\lVert#1\right\rVert}
\newcommand{\pair}[2]{\langle #1,#2\rangle}
\newcommand{\A}{\mathcal A}
\newcommand{\B}{\mathcal B}
\newcommand{\code}[1]{{\small\nolinkurl{#1}}}
\title{Quartic normal equations: complete-space solvability,\protect\\
sharp inverse stability, and atomic obstructions}
\author{Project Stardust}
\date{Reviewed editorial draft, October 29, 2026}
\begin{document}
\maketitle
\begin{abstract}
Let $H$ be a real Banach space and let $\delta:H\to L^4(\mu)$ be bounded
linear with $c\norm{u}\leq\norm{\delta u}_4$ for some $c>0$.
The quartic normal map $\A(u)(v)=\int(\delta u)^3\delta v\,d\mu$
is a bijection from $H$ onto its continuous dual. We give a direct proof
using norm-Cauchy minimizing sequences and obtain a global inverse estimate
with exponent $1/3$. A weighted lower bound at a source gives a stronger
linear stability estimate relative to that source. An actual infinite Gaussian
first-chaos embedding has a radial normal map and regular inverse at every
nonzero datum; its sharp fixed-source estimate is derived explicitly. On infinite Poisson
$L^4$, however, the exponent $1/3$ is locally optimal at every dual datum,
including data whose inverse is nonzero almost everywhere. Every fixed
finite-dimensional set of source directions can still admit linear
stability at such a source. The results are reconstructed from a frozen
Lean development; the existence and uniform-convexity mechanisms are
classical. The purpose is a complete account of the estimates, their
hypotheses, and the distinction between finite-direction and ambient
regularity. Priority for the combined presentation remains unresolved.
\end{abstract}

\section{The equation and its regularity questions}
Fix a measure space $(\Omega,\Sigma,\mu)$. No finiteness or
$\sigma$-finiteness assumption is imposed unless explicitly stated.
All $L^p$ spaces in this paper consist of equivalence classes of real
measurable functions modulo equality almost everywhere. Let $H$ be a
real Banach space with norm $\norm{\cdot}$ and continuous dual $H^*$,
equipped with its operator norm. Suppose that $\delta:H\to L^4(\mu)$
is bounded and linear and that
\begin{equation}\label{eq:lower}
 c\norm{u}\leq\norm{\delta u}_4\qquad(u\in H),\qquad c>0.
\end{equation}
Write $J(F)=F^3\in L^{4/3}(\mu)$ and define
\begin{equation}\label{eq:normal}
 \A(u)(v)=\int_\Omega (\delta u)^3\delta v\,d\mu,
 \qquad u,v\in H.
\end{equation}
The notation $\delta^*J(\delta u)$ is permissible if $\delta^*$ means
the transpose of $\delta$ applied to the integral functional represented
by $J(\delta u)$. We do not identify all of $(L^4)^*$ with $L^{4/3}$
as a preliminary assumption. This matters on general measure spaces.

The lower bound in \eqref{eq:lower} is quantitative. Injectivity alone
does not replace it in infinite dimension. It makes
$u\mapsto\norm{\delta u}_4$ an equivalent norm and makes the image of
$\delta$ closed. Indeed, if $\delta u_n$ converges in $L^4$, the lower
bound makes $u_n$ Cauchy, and completeness and continuity identify the
limit. This elementary observation explains why the source norm is
controlled throughout the paper.

\begin{theorem}\label{thm:overview}
Under \eqref{eq:lower}, the map $\A:H\to H^*$ is a smooth bijection.
Its inverse $T$ satisfies, for all $f,g\in H^*$,
\begin{equation}\label{eq:holder}
 \norm{Tf-Tg}^3\leq\frac4{c^4}\norm{f-g}.
\end{equation}
If $H$ contains a nonzero vector, no exponent $q>1/3$ gives a local
H\"older estimate for $T$ at zero. If $H=L^4(\mu_\lambda)$,
$\delta$ is the identity, and $\mu_\lambda$ is the Poisson law of a
parameter $\lambda>0$, no such exponent gives a local estimate at
\emph{any} datum of $H^*$.
\end{theorem}

The final assertion is stronger than failure of a uniform global
Lipschitz constant. It excludes even an estimate comparing one fixed
datum to nearby data. Conversely, we will prove such a fixed-datum
linear estimate under an additional weighted lower bound. Its constant
need not work between two arbitrary nearby data.

The existence argument belongs to the classical geometry of uniformly
convex spaces, initiated by Clarkson~\cite{clarkson}. Monotone-operator
solvability provides another classical context~\cite{minty,browder}.
We do not claim a new general surjectivity principle or a new
uniform-convexity inequality. The direct proof is useful here because
it exposes the precise quantitative input and avoids a separate
reflexivity construction. In classical analysis the closed embedding
in $L^4$ already implies reflexivity; avoiding that theorem in the proof
does not produce nonreflexive examples satisfying \eqref{eq:lower}.

\section{The actual duality map and its energy}
H\"older's inequality gives
\begin{equation}\label{eq:Jn}
 \norm{J(F)}_{4/3}=\norm{F}_4^3,
 \qquad \int J(F)F\,d\mu=\norm{F}_4^4.
\end{equation}
For example, $\int|F^3|^{4/3}=\int|F|^4<\infty$, which proves
membership and the norm identity directly. Almost-everywhere equality
of two representatives implies the same equality for their cubes,
so $J$ is well defined on the quotient space. The product
$J(F)G$ is integrable for every $G\in L^4$.
Consequently \eqref{eq:normal} defines a continuous functional, with
\begin{equation}\label{eq:Anorm}
 \norm{\A(u)}\leq\norm{\delta}\norm{\delta u}_4^3,
 \qquad \pair{\A(u)}u=\norm{\delta u}_4^4.
\end{equation}

Multiplication is a bounded bilinear map $L^4\times L^4\to L^2$
and also $L^2\times L^4\to L^{4/3}$, both of norm at most one.
The cube map is the composition of these maps along the diagonal.
It is therefore a continuous cubic polynomial, in particular smooth.
Composition with $\delta$ and the bounded transpose pairing proves
smoothness of $\A$. This argument concerns Fr\'echet differentiability
in the specified Banach norms; pointwise differentiation alone would
not establish it.

For $f\in H^*$ put
\begin{equation}\label{eq:energy}
 E_f(u)=\frac14\norm{\delta u}_4^4-f(u).
\end{equation}
\begin{lemma}\label{lem:derivative}
The energy is continuously Fr\'echet differentiable, and
$DE_f(u)=\A(u)-f$.
\end{lemma}
\begin{proof}
Write $F=\delta u$ and $G=\delta h$. Expansion of the quartic yields
\[
 \frac14\int(F+G)^4-\frac14\int F^4-\int F^3G
 =\frac32\int F^2G^2+\int FG^3+\frac14\int G^4.
\]
Every integral is absolutely convergent by H\"older's inequality.
The absolute value of the right side is at most
\[
 \frac32\norm F_4^2\norm G_4^2
 +\norm F_4\norm G_4^3+\frac14\norm G_4^4.
\]
Since $\norm G_4\leq\norm\delta\norm h$, this is $o(\norm h)$.
The linear term is exactly $\A(u)(h)$. Smoothness of $\A$, established
above, gives continuity of the derivative.
\end{proof}

\begin{lemma}\label{lem:below}
The range of $E_f$ is bounded below. More precisely, with
$R=\max\{1,4\norm f/c^4\}$, one has $E_f(u)\geq-\norm f R$.
\end{lemma}
\begin{proof}
If $E_f(u)\leq0$ and $s=\norm u$, then
$c^4s^4\leq4\norm f s$. For $s>1$ this implies
$c^4s\leq4\norm f$, hence $s\leq R$; for $s\leq1$ the bound
is immediate. Thus $E_f(u)\geq-f(u)\geq-\norm f R$ on the zero
sublevel. Outside that sublevel the same lower bound holds because
the energy is positive and $\norm f R\geq0$.
\end{proof}

\section{Minimization without weak compactness}
The following elementary lemma isolates the whole existence mechanism.
It will be applied after the required deficit has been proved for the
actual quartic integral.

\begin{lemma}[Quartic midpoint minimization]\label{lem:min}
Let $X$ be a real Banach space and let $E:X\to\R$ be continuous and
bounded below. If $k>0$ and
\begin{equation}\label{eq:midabstract}
 E\left(\frac{x+y}2\right)+k\norm{x-y}^4
 \leq\frac{E(x)+E(y)}2\qquad(x,y\in X),
\end{equation}
then $E$ attains its infimum.
\end{lemma}
\begin{proof}
Let $\alpha=\inf E(X)$, a finite real number. Choose $x_n$ with
$E(x_n)\to\alpha$. Since $\alpha\leq E((x_n+x_m)/2)$,
\eqref{eq:midabstract} implies
\[
 k\norm{x_n-x_m}^4\leq
 \frac{E(x_n)-\alpha+E(x_m)-\alpha}{2}.
\]
Given $\varepsilon>0$, eventually both energy differences are less
than $k\varepsilon^4$. The displayed inequality then gives
$\norm{x_n-x_m}<\varepsilon$. Completeness supplies a norm limit
$x$, and continuity gives $E(x)=\alpha$.
\end{proof}

\begin{lemma}\label{lem:mid}
The energy \eqref{eq:energy} satisfies
\begin{equation}\label{eq:mid}
 E_f\left(\frac{u+v}2\right)+\frac{c^4}{64}\norm{u-v}^4
 \leq\frac{E_f(u)+E_f(v)}2.
\end{equation}
\end{lemma}
\begin{proof}
For real $a,b$, direct expansion gives
\[
 \frac{a^4+b^4}{2}-\left(\frac{a+b}2\right)^4
 -\frac{(a-b)^4}{16}=\frac38(a^2-b^2)^2\geq0.
\]
Substitute representatives of $\delta u$ and $\delta v$ and
integrate. All fourth powers are integrable, so the resulting
inequality is valid for the quotient classes. Multiplication by $1/4$
gives the deficit $\norm{\delta(u-v)}_4^4/64$ for the potential.
The forcing term cancels exactly at the midpoint by linearity of $f$.
Finally apply \eqref{eq:lower} to $u-v$.
\end{proof}

\begin{proposition}\label{prop:onto}
For each $f\in H^*$ there exists $u\in H$ with $\A(u)=f$.
\end{proposition}
\begin{proof}
Lemmas~\ref{lem:below}, \ref{lem:mid}, and \ref{lem:min} produce a
global minimum $u$ of $E_f$. For any $h\in H$, the real function
$s\mapsto E_f(u+sh)$ has a minimum at zero, so its derivative there
vanishes. Lemma~\ref{lem:derivative} identifies that derivative with
$(\A(u)-f)(h)$. Since this holds for every $h$, $\A(u)=f$.
\end{proof}

The proof uses neither compactness of a norm ball nor an unproved
weakly convergent subsequence. The distinction is useful in a formal
development: the only limiting step is completeness applied to a
sequence already proved Cauchy. The argument also covers the zero
space, where the equation has its unique trivial solution.

\section{Uniform monotonicity and global inverse stability}
Existence did not need a uniqueness assumption. The same scalar
quartic geometry supplies uniqueness and the precise modulus of the
inverse.

\begin{lemma}\label{lem:mono}
For $u,v\in H$,
\begin{equation}\label{eq:mono}
 \pair{\A(u)-\A(v)}{u-v}
 \geq\frac14\norm{\delta(u-v)}_4^4
 \geq\frac{c^4}{4}\norm{u-v}^4.
\end{equation}
\end{lemma}
\begin{proof}
The scalar identity
\[
 (a^3-b^3)(a-b)-\frac14(a-b)^4
 =\frac34(a-b)^2(a+b)^2
\]
has a nonnegative right side. Its terms are integrable after substituting
$a=\delta u$, $b=\delta v$, by the $L^{4/3}$--$L^4$ pairing and
quartic integrability. Integration gives the first inequality, and
\eqref{eq:lower} gives the second.
\end{proof}

\begin{proof}[Proof of bijectivity and \eqref{eq:holder}]
If $\A(u)=\A(v)$, \eqref{eq:mono} forces $u=v$.
Together with Proposition~\ref{prop:onto}, this defines the inverse
$T:H^*\to H$. For $u=Tf$, $v=Tg$, dual norm control gives
\[
 \frac{c^4}{4}\norm{u-v}^4
 \leq(f-g)(u-v)\leq\norm{f-g}\norm{u-v}.
\]
When $u\neq v$, divide by $\norm{u-v}$; otherwise the desired estimate
is immediate. Taking cube roots gives the global H\"older bound
\[
 \norm{Tf-Tg}\leq(4/c^4)^{1/3}\norm{f-g}^{1/3}.
\]
\end{proof}

\begin{proposition}\label{prop:origin}
Suppose $u\neq0$. For every $q>1/3$, $K\geq0$, and $r>0$, there
exists $f\in H^*$ with $\norm f<r$ and
\[
 \norm{Tf-T0}>K\norm f^q.
\]
\end{proposition}
\begin{proof}
The defining integral gives $\A(au)=a^3\A(u)$ for real $a$ and
$\A(0)=0$. Put $f_a=\A(au)$ for $a>0$. Then
$Tf_a=au$, $T0=0$, and $\norm{f_a}=a^3\norm{\A(u)}\to0$.
The proposed estimate would require
\[
 \norm u\leq K\norm{\A(u)}^q a^{3q-1}.
\]
The right side tends to zero, whereas $\norm u>0$. Choosing $a$
small enough gives the strict reverse inequality and $\norm{f_a}<r$.
\end{proof}

The exponent is optimal in this precise local sense; no claim that
the global constant is optimal for every $\delta$ is needed. The
singularity at zero follows from cubic homogeneity. To understand
nonzero data we need information about multiplication by the source.

\section{Weighted source control and finite-dimensional directions}
For $u\in H$, define the bounded linear map
\begin{equation}\label{eq:weighted}
 \B_u:H\to L^2(\mu),\qquad \B_uw=(\delta u)(\delta w).
\end{equation}
H\"older's inequality gives
$\norm{\B_uw}_2\leq\norm{\delta u}_4\norm\delta\norm w$.
This is an actual $L^2$ product, with equality of representatives
understood almost everywhere.

\begin{lemma}\label{lem:weighted}
For all $u,v\in H$,
\begin{equation}\label{eq:wmono}
 \pair{\A(u)-\A(v)}{u-v}\geq\frac34\norm{\B_u(u-v)}_2^2.
\end{equation}
\end{lemma}
\begin{proof}
For real $a,b$,
\[
 (a^3-b^3)(a-b)-\frac34a^2(a-b)^2
 =(a-b)^2(b+a/2)^2\geq0.
\]
Integration gives the assertion. The square of the product is
integrable because $(\delta u)(\delta(u-v))\in L^2$.
\end{proof}

\begin{theorem}[Stability relative to one datum]\label{thm:weighted}
If $\beta>0$ and $\norm{\B_uw}_2\geq\beta\norm w$ for every $w\in H$,
then, for every $v\in H$,
\begin{equation}\label{eq:linear}
 \norm{u-v}\leq\frac4{3\beta^2}\norm{\A(u)-\A(v)}.
\end{equation}
Under \eqref{eq:lower}, setting $f_0=\A(u)$ yields the corresponding
bound $\norm{Tf_0-Tf}\leq4\norm{f_0-f}/(3\beta^2)$ for all $f\in H^*$.
\end{theorem}
\begin{proof}
By Lemma~\ref{lem:weighted} and the dual norm inequality,
\[
 \frac34\beta^2\norm{u-v}^2
 \leq\norm{\A(u)-\A(v)}\norm{u-v}.
\]
Divide by the distance when it is positive. The zero case is immediate.
\end{proof}

\begin{proposition}\label{prop:finite}
Suppose $\delta$ is injective and $\delta u$ is nonzero almost everywhere.
For each linear subspace $E\subset H$ of finite dimension, there is
$K_E>0$ such that
\begin{equation}\label{eq:finite}
 \norm{u-v}\leq K_E\norm{\A(u)-\A(v)}
 \qquad\text{whenever }u-v\in E.
\end{equation}
In particular, if $H$ is finite dimensional the estimate holds for
every $v\in H$.
\end{proposition}
\begin{proof}
If $\B_uw=0$, then $(\delta u)(\delta w)=0$ almost everywhere.
Nonvanishing of $\delta u$ gives $\delta w=0$, and injectivity gives
$w=0$. Thus $\B_u$ is injective. Restricted to $E$, it has a positive
lower norm bound: for $E\neq\{0\}$, minimize the continuous positive
function $w\mapsto\norm{\B_uw}_2$ on the compact unit sphere of $E$.
Its minimum $\beta_E$ is positive; scaling gives the bound on $E$.
For $E=\{0\}$ any positive lower-bound constant works. Apply the proof
of Theorem~\ref{thm:weighted} to $w=u-v\in E$.
\end{proof}

The finite-dimensional constant depends on the subspace. The
proposition provides no uniform lower bound as $E$ varies. This
dependence is decisive on an infinite atomic space.

As a regular example, let $\mu$ be a probability measure, $H=\R$,
and $\delta u$ the constant function $u$. Then
$\norm{\delta u}_4=|u|$ and $\norm{\B_uw}_2=|u||w|$.
At $u\neq0$, \eqref{eq:linear} gives a linear stability constant
$4/(3|u|^2)$. Under the usual identification of $\R^*$ with $\R$,
the normal equation is $u^3=f$. The behavior at zero and at nonzero
data is thus already different in one dimension.


\section{Gaussian synthesis and inverse regularity away from zero}
The infinite-dimensional failure below is not a consequence of dimension
alone. Consider the actual Gaussian product space carrying independent
standard coordinates $(G_j)_{j\geq0}$, and take $H=\ell^2(\N;\R)$.
Its isonormal embedding $W(h)$ is the $L^2$ limit of
$\sum_{j<N}h_jG_j$. Every finite linear combination of such limits is
Gaussian: convergence in $L^2$ implies convergence of characteristic
functions, whose finite formulas tend to
$\exp(-\|h\|^2t^2/2)$. In particular
\[
 \mathbb E W(h)^4=3\|h\|^4,
 \qquad \delta h:=W(h)\in L^4,
 \qquad \|\delta h\|_4=3^{1/4}\|h\|.
\]
Linearity holds on the actual equivalence classes: it first holds in
$L^2$, hence almost everywhere, and the representatives belong to
$L^4$. Thus $\delta:H\to L^4$ is a bounded closed embedding satisfying
\eqref{eq:lower}. It is not the identity on the whole Gaussian $L^4$.

For $h,k\in H$, the joint moment generating function is
$\exp((\|h\|^2s^2+2\langle h,k\rangle st+\|k\|^2t^2)/2)$.
Its $s^2t^2$ and $s^3t$ coefficients give
\begin{align}\label{eq:gaussianweighted}
 \|\B_hk\|_2^2
 &=\mathbb E[W(h)^2W(k)^2]
   =\|h\|^2\|k\|^2+2\langle h,k\rangle^2,\\
 \A(h)(k)&=\mathbb E[W(h)^3W(k)]
   =3\|h\|^2\langle h,k\rangle.\nonumber
\end{align}
These are genuine finite moments by Gaussian integrability and
H\"older's inequality, including degenerate source pairs.
Hence $\|\B_hk\|_2\geq\|h\|\|k\|$ on the \emph{entire}
infinite source space. At every $h\ne0$, Theorem~\ref{thm:weighted}
already gives the mapped fixed-datum estimate
\begin{equation}\label{eq:gaussianmapped}
 \|h-Tf\|\leq\frac4{3\|h\|^2}\|\A(h)-f\|.
\end{equation}
The source \code{GaussianFirstChaosQuarticSource} proves the actual
$L^4$ realization, these two moment identities and this estimate.

\subsection{Exact radial inverse and a sharp fixed-datum constant}
Identify $H^*$ with $H$ by the real Riesz isometry
$Rh(k)=\langle h,k\rangle$. Equation~\eqref{eq:gaussianweighted}
identifies the normal map with
\[
 N(h)=3\|h\|^2h.
\]
Its inverse, in these Hilbert coordinates, is
\begin{equation}\label{eq:gaussianinverse}
 T(y)=\begin{cases}
 3^{-1/3}\|y\|^{-2/3}y,&y\ne0,\\0,&y=0.
 \end{cases}
\end{equation}
Indeed $\|N(h)\|=3\|h\|^3$, and substitution checks both inverse
identities. In particular the origin retains the cubic-root singularity.
The following sharper bound and its optimality are explicit human
consequences of the mapped radial identity, not additional named
formal declarations.

\begin{proposition}\label{prop:gaussiansharp}
For every fixed $h\ne0$ and every $k\in H$,
\begin{equation}\label{eq:gaussiansharp}
 \|h-k\|\leq\frac4{9\|h\|^2}\|N(h)-N(k)\|.
\end{equation}
The constant is best possible for comparison with this fixed source
against all other sources, already on its one-dimensional span.
Equivalently it bounds $\|T(N(h))-T(y)\|$ for every datum $y$.
\end{proposition}
\begin{proof}
Put $r=\|h\|$, $s=\|k\|$, $c=\langle h,k\rangle$ and
$d=r^2-rs+s^2$. Direct expansion of squared Hilbert norms gives
\begin{align*}
 &\|r^2h-s^2k\|^2-d^2\|h-k\|^2\\
 &\hspace{12mm}=2(r^2+s^2)(r-s)^2(c+rs)\geq0.
\end{align*}
The sign uses Cauchy--Schwarz, $c\geq-rs$. Both norms and $d$
are nonnegative, so taking square roots proves
$\|r^2h-s^2k\|\geq d\|h-k\|$.
Moreover $d=(s-r/2)^2+3r^2/4\geq3r^2/4$.
Multiplication by three gives \eqref{eq:gaussiansharp}.
For $k=-h/2$, the source distance is $3r/2$ and the image distance
is $27r^3/8$. Their ratio is $4/(9r^2)$, proving optimality.
\end{proof}

\subsection{Local derivatives and the distinction between estimates}
Away from zero, the scalar norm powers in \eqref{eq:gaussianinverse}
are continuously differentiable. Since
$D(\|y\|^2)[v]=2\langle y,v\rangle$, the scalar power rule and
product rule give, at $y=N(h)$,
\begin{equation}\label{eq:gaussianderivative}
 DT(y)=\frac1{3\|h\|^2}\left(I-\frac23\Pi_h\right),
 \qquad \Pi_hv=\frac{\langle h,v\rangle}{\|h\|^2}h.
\end{equation}
The derivative acts by $1/(9\|h\|^2)$ along $h$, and by
$1/(3\|h\|^2)$ on its orthogonal complement. In this infinite
Hilbert space the latter complement is nonzero, so the derivative
norm is exactly $1/(3\|h\|^2)$. Formula~\eqref{eq:gaussiansharp}
is a different, global comparison relative to a fixed nonzero datum;
its sharp constant is larger than this local derivative norm.

For completeness, the inverse also has ordinary two-point local
Lipschitz regularity. On a datum ball
$\|y-y_0\|<\rho<\|y_0\|$, its derivative norm is bounded by
$3^{-1/3}(\|y_0\|-\rho)^{-2/3}$. The segment joining two points
of this ball stays in the ball. Integrating the continuous derivative
along that segment gives the same bound times their distance.
This is an elementary human differentiation of the explicit inverse,
not an inference that the general one-base estimate alone implies
an arbitrary two-point estimate.

The underlying radial map is the classical $p=4$ vector power map;
Lindqvist~\cite{lindqvist}, Section 10, gives the standard vector
power identities and monotonicity framework. The exact fixed-source
constant above is proved directly here; the bounded comparison does
not establish its priority in the literature. The contrast with the
next section changes the source embedding: Gaussian first-chaos
synthesis on $\ell^2$ has uniform weighted control, whereas the full
Poisson $L^4$ identity permits concentrated atomic perturbations.
Neither infinite dimension nor nonvanishing of an individual source
settles inverse regularity without that geometry.

\section{Infinite Poisson space and failure of weighted coercivity}
Fix $\lambda>0$. Give $\N=\{0,1,\ldots\}$ its full power-set
$\sigma$-algebra and Poisson probability measure
\[
 p_n=\mu_\lambda(\{n\})=e^{-\lambda}\frac{\lambda^n}{n!}>0.
\]
In this section $H=L^4(\mu_\lambda)$ and $\delta$ is the identity.
The lower bound \eqref{eq:lower} holds with $c=1$. Every coordinate
has positive mass, so an equivalence class has well-defined values
$F(n)$. Its norm is
\[
 \norm F_4^4=\sum_{n=0}^\infty p_n|F(n)|^4.
\]
Let $e_n=p_n^{-1/4}1_{\{n\}}$. These are normalized in $L^4$.
For a fixed $F\in H$ put $x_n=F(n)p_n^{1/4}$. Summability gives
\begin{equation}\label{eq:x}
 \sum_n|x_n|^4<\infty,\qquad x_n\longrightarrow0.
\end{equation}

\begin{proposition}\label{prop:noweight}
No $F\in L^4(\mu_\lambda)$ satisfies a positive lower bound
$\norm{Fw}_2\geq\beta\norm w_4$ for all $w\in L^4(\mu_\lambda)$.
\end{proposition}
\begin{proof}
Taking $w=e_n$ gives
\[
 \norm{Fe_n}_2=|F(n)|p_n^{1/4}=|x_n|\longrightarrow0,
 \qquad\norm{e_n}_4=1.
\]
A positive lower bound contradicts this limit.
\end{proof}

Even the constant source $F=1$ has this obstruction. It is nonzero
everywhere, and multiplication by $F$ is injective, but the inclusion
of this probability-space $L^4$ into $L^2$ has no positive lower bound.
This does not yet prove failure of inverse Lipschitz stability:
Theorem~\ref{thm:weighted} is a sufficient criterion, not an equivalence.
We next prove the stronger failure directly by perturbing the normal
map itself.

\section{The inverse exponent is optimal at every atomic datum}
The integral pairing represented by $e_n^3$ has dual norm at most one
by H\"older's inequality. For $a\geq0$, the functions $F+ae_n$ and
$F$ differ at only one coordinate. Cubing gives the exact identity
\begin{equation}\label{eq:atom}
 \A(F+ae_n)-\A(F)
 =\big((x_n+a)^3-x_n^3\big)\A(e_n).
\end{equation}
In particular,
\begin{equation}\label{eq:shift}
 \norm{\A(F+ae_n)-\A(F)}
 \leq a\big(3|x_n|^2+3a|x_n|+a^2\big).
\end{equation}
To check \eqref{eq:atom} without confusing values and normalized
coordinates, observe that on the $n$th atom
$F(n)=p_n^{-1/4}x_n$ and $e_n(n)=p_n^{-1/4}$.
The cube difference is therefore $p_n^{-3/4}((x_n+a)^3-x_n^3)$,
exactly the multiplier of $e_n^3$. Off that atom both sides vanish.

\begin{theorem}\label{thm:everywhere}
For every $f_0\in H^*$, $q>1/3$, $K\geq0$, and $r>0$, there exists
$f\in H^*$ such that
\begin{equation}\label{eq:everywhere}
 \norm{f-f_0}<r,\qquad
 \norm{Tf-Tf_0}>K\norm{f-f_0}^q.
\end{equation}
\end{theorem}
\begin{proof}
By surjectivity write $F=Tf_0$. Let
\[
 a_n=|x_n|+\frac1{n+1},\qquad U_n=F+a_ne_n,\qquad f_n=\A(U_n).
\]
Then $a_n>0$ and $a_n\to0$ by \eqref{eq:x}. Since $|x_n|\leq a_n$,
\eqref{eq:shift} gives
\[
 \norm{f_n-f_0}\leq7a_n^3\longrightarrow0.
\]
Injectivity of $\A$ and $\norm{e_n}_4=1$ give
$\norm{Tf_n-Tf_0}=\norm{U_n-F}_4=a_n$.
Because $3q-1>0$, eventually
$K7^q a_n^{3q-1}<1$. Hence
\[
 K\norm{f_n-f_0}^q\leq K7^q a_n^{3q}<a_n.
\]
Choose $n$ large enough for this inequality and the radius condition.
\end{proof}

This completes the optimality assertion in Theorem~\ref{thm:overview}.
The proof uses small normalized coordinates of an arbitrary $L^4$
source, rather than a zero coordinate. Thus it applies to strictly
positive sources as well as sources supported on finitely many atoms.
The Poisson distribution is a concrete verified model; the argument
itself only uses the displayed positive atomic weights and quartic
summability. We do not claim a separate Lean theorem for every atomic
measure under that generalization.

\begin{corollary}\label{cor:directions}
At the constant source $F=1$, every fixed finite-dimensional source
subspace $E\subset L^4(\mu_\lambda)$ admits \eqref{eq:finite}. There
is no common constant valid for all one-dimensional source subspaces.
\end{corollary}
\begin{proof}
Apply Proposition~\ref{prop:finite} to $\delta=\mathrm{id}$ and $F=1$.
If a constant $K$ worked for all one-dimensional subspaces, then for
each $V$ take $E=\operatorname{span}\{1-V\}$. The estimate would hold
for every $V$. Taking $V=Tf$ contradicts Theorem~\ref{thm:everywhere}
with $f_0=\A(1)$ and $q=1$.
\end{proof}

The order of quantifiers is the issue: for each fixed $E$ there is a
constant, but no single constant works for all $E$. A finite-dimensional
calculation at a nonvanishing source can therefore be accurate while
providing no uniform ambient regularity theorem.

\section{Disjoint sources and the role of support}
A simpler obstruction at selected nonzero data helps separate two
mechanisms. Assume \eqref{eq:lower}. Suppose $u,v\in H$, $v\neq0$,
and $(\delta u)(\delta v)=0$ almost everywhere. Then mixed terms in
the cubic vanish, so
\begin{equation}\label{eq:disjoint}
 \A(u+av)=\A(u)+a^3\A(v)\qquad(a\in\R).
\end{equation}
\begin{proposition}\label{prop:disjoint}
For every $K\geq0$ and $r>0$ there exists $f$ with
$\norm{f-\A(u)}<r$ and
$\norm{Tf-u}>K\norm{f-\A(u)}$.
\end{proposition}
\begin{proof}
Use $f_a=\A(u+av)$ for $a>0$. The source difference has norm
$a\norm v$, and \eqref{eq:disjoint} gives dual difference
$a^3\norm{\A(v)}$. Small $a$ meets both required inequalities.
\end{proof}

For the Poisson identity source one may take indicators of two
distinct atoms. This produces a nonzero datum with no linear inverse
estimate using disjoint support alone. Theorem~\ref{thm:everywhere}
goes further: it works when the source never vanishes and no nonzero
disjoint source exists. Small normalized tail coordinates, rather
than exact support separation, then replace cubic decoupling.

\section{Formal scope and reproducibility}\label{sec:formal}
The mathematical baseline is the frozen Git commit
\begin{center}\small\ttfamily
bab688bb6140cd6316c80680095a914ea0d346db.
\end{center}
All declarations below have prefix
\code{OperatorDerivativeRewrite.Foundation}. The accompanying
\code{theorem-map.json} records exact files, declarations, and hashes.
The principal correspondence is as follows.

\begin{longtable}{@{}p{0.27\textwidth}p{0.67\textwidth}@{}}
\textit{Paper result}&\textit{Lean namespace and declaration}\\\hline
\eqref{eq:Jn}, \eqref{eq:normal}&
\code{BanachLpDualityFour.dualityFour_norm},
\code{dualityFour_pairing_eq_norm}, \code{normalFour}\\[3pt]
Lemma~\ref{lem:derivative}&
\code{BanachLpDualityFourExistence.hasFDerivAt_quarticPotential}\\[3pt]
Lemmas~\ref{lem:below}--\ref{lem:mid}&
\code{BanachLpDualityFourCoercivity.quarticPotential_lower_bound},
\code{quarticEnergy_midpoint_lower};
\code{QuarticUniformMinimization.exists_minimizer_of_quartic_midpoint}\\[3pt]
Bijectivity, \eqref{eq:holder}&
\code{BanachLpDualityFourInfinite.normalFour_existsUnique},
\code{normalFour_inverse_holder}\\[3pt]
Proposition~\ref{prop:origin}&
\code{BanachQuarticHolderSharpness.no_local_holder_bound}\\[3pt]
Theorem~\ref{thm:weighted}&
\code{BanachQuarticWeightedStability.inverse_lipschitz_at_weighted_datum}\\[3pt]
Proposition~\ref{prop:finite}&
\code{BanachFiniteDirectionStability.finite_direction_stability}\\[3pt]
Proposition~\ref{prop:noweight}&
\code{PoissonQuarticWeightedObstruction.no_positive_weighted_lower}\\[3pt]
\eqref{eq:atom}, \eqref{eq:shift}&
\code{PoissonQuarticAtomPerturbation.normal_atom_shift},
\code{normal_shift_norm_bound}\\[3pt]
Theorem~\ref{thm:everywhere}&
\code{PoissonQuarticHolderEverywhere.no_local_holder_at_datum}\\[3pt]
Corollary~\ref{cor:directions}&
\code{BanachFiniteDirectionStability.poisson_constant_finite_directions},
\code{no_uniform_one_direction_bound}\\[3pt]
Proposition~\ref{prop:disjoint}&
\code{BanachQuarticDisjointObstruction.no_local_linear_bound_at_disjoint_image}\\
\end{longtable}

The complete-space namespace is essential in the existence row:
an earlier namespace contains only a finite-dimensional existence
theorem, although its energy derivative is dimension independent.
The map called \code{normalFourEquiv} is a nonlinear bijection,
not a linear equivalence or an identification supplied by a Hilbert
inner product. The paper's $H^*$ is the continuous real dual.

The baseline audit receipt records a successful complete root build,
54,943 audited declarations, and acceptance only of
\code{propext}, \code{Classical.choice}, and \code{Quot.sound}.
The editorial work changes no Lean source or dependency configuration.
It reuses that existing verification rather than claiming a new full
build. The source hashes and receipt provenance are included with the
manuscript. Lean is pinned to version 4.34.1 and Mathlib to revision
\code{d13f23b723b8a846827a245b89c10fc7d3f11612}.

\section{What the estimates do and do not decide}
The lower bound on $\delta$ gives existence on the whole continuous
dual, uniqueness, and a global modulus. A weighted lower bound gives
stronger control relative to one source. Nonvanishing almost everywhere
implies injectivity of the weighted map; finite dimension converts
that injectivity into a lower bound. Infinite Poisson space shows
exactly where that last implication fails, and direct perturbations
show that the failure corresponds to an actual inverse singularity.

No stochastic integral is constructed by these arguments. In
particular, a Banach-space lower bound for a proposed stochastic
divergence cannot be presumed from the existence theorem: it is an
input that must be established for the actual operator. The results
also do not settle general exponents, unbounded transpose domains,
or closability characterizations. These are separate questions.

The solvability and uniform monotonicity belong to established
nonlinear analysis. The bounded preliminary literature check places
them in that context and does not certify priority for the particular
combination of atomic optimality and finite-direction regularity.
This manuscript is a reviewed exposition of verified mathematics,
with mathematical novelty explicitly unresolved. An external
submission would require a more targeted comparison with the
literature on inverse duality mappings and their local moduli.

\begin{thebibliography}{9}
\bibitem{clarkson} J. A. Clarkson,
Uniformly convex spaces, \emph{Trans. Amer. Math. Soc.} \textbf{40}
(1936), 396--414. \href{https://doi.org/10.1090/S0002-9947-1936-1501880-4}{doi:10.1090/S0002-9947-1936-1501880-4}.
\bibitem{minty} G. J. Minty,
Monotone (nonlinear) operators in Hilbert space,
\emph{Duke Math. J.} \textbf{29} (1962), 341--346.
\href{https://doi.org/10.1215/S0012-7094-62-02933-2}{doi:10.1215/S0012-7094-62-02933-2}.
\bibitem{browder} F. E. Browder,
Nonlinear elliptic boundary value problems,
\emph{Bull. Amer. Math. Soc.} \textbf{69} (1963), 862--874.
\href{https://doi.org/10.1090/S0002-9904-1963-11068-X}{doi:10.1090/S0002-9904-1963-11068-X}.
\bibitem{lindqvist} P. Lindqvist,
\emph{Notes on the p-Laplace equation}, author-hosted primary notes,
Section 10, inequalities for vectors, pages 71--73.
\url{https://lqvist.folk.ntnu.no/p-laplace.pdf}.
\end{thebibliography}
\end{document}
